\documentclass[11pt]{article}
\usepackage[margin=0.75in]{geometry}
\usepackage{fontspec}
\usepackage{setspace}
\usepackage{hyperref}
\hypersetup{colorlinks=true, urlcolor=blue, linkcolor=black}
\pagestyle{empty}
\setlength{\parskip}{0.7em}
\setlength{\parindent}{0pt}
\setstretch{0.94}
\setmainfont{Times New Roman}

\begin{document}

\begin{flushleft}
Ke Cheng HSU (Clio) \\
Department of Economics, The Hong Kong University of Science and Technology \\
Clear Water Bay, Kowloon, Hong Kong $\cdot$ +852 6840 9753 \\
\href{mailto:khsu@connect.ust.hk}{khsu@connect.ust.hk} $\cdot$ \href{https://cliohsu1997.github.io/}{cliohsu1997.github.io}
\end{flushleft}

18 November 2026

\vspace{0.12em}
Professor Jr-Yan Wang \\
Chairperson, Department of International Business \\
School of Management, National Taiwan University \\
No.\ 1, Sec.\ 4, Roosevelt Road \\
Taipei 106, Taiwan

\vspace{0.12em}
Dear Professor Wang,

I am writing to apply for a full-time tenure-track faculty position (Assistant Professor) in the Department of International Business, to begin Fall 2027. I am a Ph.D.\ candidate in Economics at the Hong Kong University of Science and Technology, expecting to graduate in June 2027. My research is in \textbf{quantitative marketing} and empirical industrial organization, with a focus on digital markets, return and pricing policy, product design, and consumer search. I grew up in Taiwan, worked as a research assistant at National Taiwan University before my doctoral studies, and can teach in Mandarin and English.

My job market paper, ``Designing Return Policy in E-Commerce: A Trade-off between Purchase and Return Efficiency,'' is a quantitative-marketing study of how an online seller should design the return offer. A stricter policy curbs inefficient returns after delivery, but the lower return rate makes a higher product price more attractive for the seller and can exclude efficient purchases. Using transaction data from a Taobao apparel retailer, I estimate a structural model of purchase, return, and whether to buy return-shipping insurance under monopoly pricing across six categories. Automatic coverage of return shipping performs better when the margin, net scrap value, and uncertainty are high, because it expands efficient purchases. Optional coverage performs better when those attributes are low, because it curbs excessive returns. A menu that pairs a discounted price with a medium restocking fee balances the two margins and performs better than both insurance policies for most categories.

A second working paper, ``Product Design with Heterogeneous Search Costs'' (with Pak Hung Au), studies how firms choose product design when consumers differ in search costs. Search costs shape which products consumers search. Consumers with low search costs seek niche products, and consumers with high search costs seek standardized products. A fall in search costs can raise prices when this shift in which product consumers search outweighs the effect that encourages more search. I also have two ongoing projects. With Kohei Kawaguchi and Kei Kawai, I use Japanese bookstore data to study how resale price maintenance encourages retailers to hold inventory in low-demand states, at an environmental cost when unsold copies are returned and discarded. With Ce Bian and Liang Zhao, I study pass-through when products can be returned, and show that a tax can lower the equilibrium price and raise welfare by reducing the return rate.

At HKUST I was a teaching assistant from 2020 to 2025 for undergraduate and Ph.D.\ courses, with an average evaluation of 4.72/5, and in 2026 I designed and lectured the mathematics camp for the MSc in Economics, receiving 4.53/5. I am prepared to teach quantitative marketing, pricing and consumer analytics, digital marketing / e-commerce, managerial economics, and applied microeconometrics at the undergraduate and master's levels.

I submit my CV, research statement with current working papers, and job market paper. Crystal, the secretary of the Department of Economics at HKUST, will send three confidential recommendation letters on behalf of Pak Hung Au, Kohei Kawaguchi, and Kei Kawai to \href{mailto:stellacheng@ntu.edu.tw}{stellacheng@ntu.edu.tw}. You can reach me at \href{mailto:khsu@connect.ust.hk}{khsu@connect.ust.hk} or +852 6840 9753.

\vspace{0.25em}
Sincerely,

\vspace{0.25em}
Ke Cheng HSU \\
Ph.D.\ Candidate, Department of Economics, HKUST

\end{document}
